\subsection*{一步一步建立直观}

\paragraph{1. 把点积当作沿某个方向的读数}

把 $u$ 看成测量方向，$v$ 看成一次位移。投影回答位移沿尺子方向前进了多少；垂直于尺子的那部分不会贡献读数。若尺子没有归一化，点积还额外乘上尺子的长度，因此不能直接把它当作投影长度。



\paragraph{2. 垂直分量为什么消失}

让待测向量与测量方向垂直，点积变为零。一般情形可以把 $v$ 拆成平行部分与垂直部分；平行部分的系数必须是 $(u\cdot v)/(u\cdot u)$，因为减掉它以后与 $u$ 的点积恰为零。这里要求 $u$ 非零。



\paragraph{3. 负号表示相反的测量方向}

把待测向量移到尺子的反方向，点积为负。这个符号表达朝向，并非说长度可以为负。在力学中，恒力与位移的点积描述功；垂直运动不做功、反向运动做负功，提供了同一个投影关系的物理解释。



\subsection*{定义、命题与证明}

\paragraph{定义：沿一条非零方向的正交投影}

对非零向量 $u$，定义 $\operatorname{proj}_u(v)=\dfrac{u\cdot v}{u\cdot u}u$。投影是向量，点积是标量，两者不能互相替代。

\paragraph{命题：投影后的剩余量与投影方向垂直}

设 $u\ne0$，令 $r=v-\operatorname{proj}_u(v)$。则 $u\cdot r=0$。

\begin{proof}

由点积的线性性，$u\cdot r=u\cdot v-\dfrac{u\cdot v}{u\cdot u}(u\cdot u)=0$。分母可约去，是因为实内积下非零向量满足 $u\cdot u>0$。

因此 $v$ 被分成沿 $u$ 的部分与垂直于 $u$ 的部分。若 $u=0$，公式没有定义；图形窗口会报告这一退化情形，而不生成一个任意的投影方向。

\end{proof}
